\clearpage
\VerdierPhysicalPageAnchor{0241}%
\begingroup\small\setlength{\parskip}{0pt}\setlength{\emergencystretch}{1.5em}
\setlength{\abovedisplayskip}{4pt plus 2pt minus 2pt}%
\setlength{\belowdisplayskip}{4pt plus 2pt minus 2pt}%

\VerdierSemanticAnchor{ch4d-4-5-2}
\noindent\textbf{4.5.2.}\quad
Let $F\colon D(\mathcal A)\longrightarrow\mathcal B$ be a cohomological
functor stationary on the right, let $p$ be an integer, let $X$ be an
object of $\operatorname{comp}(\mathcal A)$, and let

\VerdierSemanticAnchor{ch4d-4-5-2-second-filtration}
\[
0=X(-\infty)\longrightarrow\cdots\longrightarrow X(q)
\longrightarrow X(q+1)\longrightarrow\cdots
\longrightarrow X(+\infty)=X
\]
be the second canonical filtration (4.3.1) of $X$. The inductive system
$F^p(X(q))$ is stationary; i.e.\ there is an integer $q_0$ such that, for
every $q$ with $q_0\leq q\leq+\infty$, the morphism
\[
F^p(X(q))\longrightarrow F^p(X)
\]
is an isomorphism. Indeed, with the notation of (4.2.3), there is an exact
sequence
\[
F^{p-1}\bigl(X(q,+\infty)\bigr)\longrightarrow F^p(X(q))
\longrightarrow F^p(X)\longrightarrow F^p\bigl(X(q,+\infty)\bigr),
\]
and, since the functor $F$ is stationary on the right, there is an integer
$q_0$ such that, for every $q>q_0$, the objects $F^p(X(q,+\infty))$ and
$F^{p-1}(X(q,+\infty))$ are zero.

\VerdierSemanticAnchor{ch4d-4-5-2-first-filtration}
Likewise, let $F\colon D(\mathcal A)\longrightarrow\mathcal B$ be a
cohomological functor stationary on the right, let $p$ be an integer, let
$Y$ be an object of $\operatorname{comp}(\mathcal A)$, and let
\[
0=Y(-\infty)\longrightarrow\cdots\longrightarrow Y(q)
\longrightarrow Y(q+1)\longrightarrow\cdots
\longrightarrow Y(+\infty)=Y,
\]
be the first filtration of $Y$ (4.1.2). There is an integer $q_0$ such that,
for every $q<q_0$, the object $F^p(Y(q))$ is zero. The verification analogous
to the preceding one is left to the reader, who may likewise pass to the
opposite categories to study the analogous properties of functors
stationary on the left.

\VerdierSemanticAnchor{ch4d-4-5-3}
\medskip\noindent\textbf{Proposition 4.5.3.}\quad

\VerdierSemanticAnchor{ch4d-4-5-3-a}
\noindent\textit{a) Let $Y$ be an object of
$\operatorname{comp}^{*}(\mathcal A)$ (1.1.1)
($*=+,-,b,{}$ ``empty''). Let
$F\colon D^*(\mathcal A)\longrightarrow\mathcal B$ be a cohomological
functor. The first spectral sequence of the functor $F$ relative to the
complex $Y$ (4.4.1), whose initial terms are (4.4.2)}
\[
\textit{$\mathrm{I}_1^{p,q}=F^q(Y^p)$}
\]
\textit{is stationary (Chapter~II, 4.4.2) in the following cases:}

\begin{enumerate}
\item[1)] \textit{$Y$ is an object of $\operatorname{comp}^{b}(\mathcal A)$.}
\item[2)] \textit{$Y$ is an object of $\operatorname{comp}^{+}(\mathcal A)$
and $F$ is stationary on the right.}
\item[3)] \textit{$Y$ is an object of $\operatorname{comp}^{-}(\mathcal A)$
and $F$ is stationary on the left.}
\item[4)] \textit{The functor $F$ is stationary.}
\end{enumerate}

\VerdierSemanticAnchor{ch4d-4-5-3-b}
\noindent\textit{b) Let $X$ be an object of $D^*(\mathcal A)$ (1.2.4), and
let $F\colon D^*(\mathcal A)\longrightarrow\mathcal B$ be a cohomological
functor. The second spectral sequence of the functor $F$ relative to the
complex $X$ (4.4.5) is stationary in the following cases:}
\endgroup
